% concatenated from Article___SNP_Weyl_Connections__1_.tex
\documentclass[12pt]{article}


\usepackage[T1]{fontenc}
\usepackage{avant}
\usepackage[cp1250]{inputenc}
\usepackage{amsmath,amssymb,amsthm}
\usepackage{amscd}
\usepackage{mathtools} % For coloneqq :=
\usepackage{stmaryrd} % Various symbols like llbracket
\usepackage[ruled,linesnumbered,norelsize]{algorithm2e} %for algorithm
\usepackage{longtable} %for page break for tables
\usepackage{url} % for hyperlinks
\usepackage{authblk} % for affiliation

\usepackage{tikz}
%\usetikzlibrary{matrix,arrows,calc}
\usepackage[all]{xy}
\newtheorem{theorem}{Theorem}
\newtheorem{proposition}{Proposition}
\newtheorem{lemma}{Lemma}
\newtheorem{definition}{Definition}
\newtheorem{example}{Example}
\newtheorem{note}{Note}
\newtheorem{corollary}{Corollary}
\newtheorem{remark}{Remark}
\newtheorem{question}{Question}
\newtheorem{problem}{Problem}
\newtheorem{answer}{Answer}
\newtheorem{conjecture}{Conjecture}
\newtheorem{comment}{Comment} 


\title{On non-positive Weyl connections on  Lie groups}
\author{Tomasz Dudek and Aleksy Tralle}



\begin{document}
\maketitle{}


\section{Introduction}
Let there be given a Riemannian manifold $(M,\langle -,-\rangle)$. Let $\nabla$ denote the Levi-Civita coonection.
By definition, a Weyl connection is given by the formula
$$\hat\nabla_XY=\nabla_XY+\varphi(Y)X+\varphi(X)Y-\langle X,Y\rangle E,$$
where $\varphi$ is a one-form $\varphi(X)=\langle X,E\rangle$, $X,Y$ are arbitrary vector ields, $E$ is fixed.
Given a Weyl connection $\widehat{\nabla}$ one considers the curvature tensor
$$\widehat{R}(X,Y)=\widehat{\nabla}_XY\widehat{\nabla}_Y-\widehat{\nabla}_Y\widehat{\nabla}_X-\widehat{\nabla}_{[X,Y]}$$ 
and takes the {\it anti-symmetric} part $\widehat{R}_a$ of the tensor $\widehat{R}$. This tensor is called {\it the directional curvature}.
\begin{definition} The Weyl sectional curvature {\rm is defined as}
$$\widehat{K}(\Pi)=\langle \widehat{R}_a(X,Y)Y,X\rangle,\Pi=
\operatorname{Span}_{\mathbb{R}}(X,Y),\,|X|=|Y|=1, X\perp Y.$$
\end{definition}
\begin{proposition}[\cite{TW}] The sign of $\widehat{K}(\Pi)$ is well defined.
\end{proposition}
There are several reasons for an interest in studying such connections and their curvature properties. First, we mention the applications in dynamical systems. In \cite{W} the notion of W-flow on the unit tangent bundle to $(M,\langle -,-\rangle)$ was introduced. The orbits of such flow are geodesics of $\widehat{\nabla}$. An interesting theorem in \cite{W}, Theorem 5.1  says that (in particular) on a Riemannian manifold with a Weyl connection whose {\it Weyl curvature} is negative everywhere the W-flow $\Phi$ has dominated splitting, that is, the tangent bundle continuously splits as 
$TM=E\oplus F$ in a way that the subbundles $E$ and $F$ are invariant with respect to the differential $\Phi_*: TM\rightarrow TM$ and  have the following property: 
$$||\Phi_*)^N_xv||\leq {1\over 2}||(\Phi_*)_x(w)||,\qquad \forall v\in E_x, 
\quad w\in F_x.$$ 
It follows that W-flows on such compact  manifolds  have a kind of "hyperbolic" properties akin to the Anosov property. 

The second reason comes from the theory of Gaussian thermostats \cite{H} which yields interesting models in statistical physics.

 It follows that a problem of  description of Riemannian manifolds with negative Weyl connections is intriguing, but difficult, and still not much is known. In particular, the following problem was posed in \cite{W}.         

\begin{problem}\label{probl-gen}
 Construct a Weyl structure $(M,\langle -,-\rangle,\widehat{\nabla})$ with negative Weyl curvature on a manifold which does not carry a Riemannian metric of negative sectional curvature.
\end{problem}
It is quite natural to begin with the case of homogeneous spaces $G/H$ of Lie groups and invariant Weyl connections, since in this case one can use the power of Lie theory. 
\begin{problem}\label{probl:inv}
Describe invariant Weyl connections on homogeneous spaces and solve Problem \ref{probl-gen} in the homogeneous case. 
\end{problem}
The first results in this direction were obtained in \cite{TW}. In this work, the authors propose a way to look for Weyl connections with {\it non-positive} Weyl curvature by rescaling the vector field $E$ which determines $\widehat{\nabla}$.   
\begin{definition} {\rm By definition a {\it stretching} of $\widehat{\nabla}$ is a one-parameter  family $\widehat{\nabla}_t$ of Weyl connections}
$$\widehat{\nabla}_t=\nabla_XY+t\varphi(X)Y+t\varphi(Y)X-t^2\langle X, Y\rangle E.$$
A connection $\widehat{\nabla}$ is called stretched non-positive (SNP) if there exists a stretching such that $\widehat{\nabla}_t$ is non-positive for sufficiently large $t$.
\end{definition}     
Since the problem (even for SNP) is difficult for arbitrary homogeneous spaces, it is natural to begin with invariant Riemannian metrics on a Lie group $G$  and consider  Lie groups endowed with left-invariant Riemannian metrics and invariant Weyl connections determined by left-invariant vector fields $E\in\mathfrak{g}$ (that is, vectors in the corresponding Lie algebra $\mathfrak{g}$). In \cite {TW} the authors analyzed the case of 3-dimensional Lie groups and found an (isolated) example of non-positive on a solvable Lie group SOL. They conjecture that in 3-dimensional case, only SOL carries such connections. We confirm this in Theorem \ref{thm:main} in the case of completely solvable Lie groups.  In dimension 4,  SNP connections were found on some 4-dimensional solvable Lie groups, and the structure of solvable Lie groups endowed with SNP connctions determined by parallel vector fields $E$ was described in \cite{BJT}. In this note we continue this line of research. In particular we prove the following. 
\begin{theorem}\label{thm:main} A 3-dimensonal unimodular completely solvable Lie group admits  a left-invariant Riemannian metric and a left- invariant Weyl connection $\widehat{\nabla}$ determined by a choice of $E\in\mathfrak{g}$ which has the SNP property if and only if it is the group SOL.
\end{theorem}
Our approach is based on the following:  we express the SNP condition in terms of the Nomizu function and simplify the calculations using the classification of the invariant Riemannian metrics on completely solvable Lie groups up to Lie algebra isomorphisms \cite{KHL}.  We believe that this approach can be extended to  arbitrary completerly solvable Lie groups.   
  
\section{Invariant connections on homogeneous spaces of Lie groups}
We will follow more modern exposition of the classical Nomizu's theorem in \cite{DGP}. In this section $\nabla$ denotes arbitrary affine connection (not necessarily Levi-Civita or Weyl). Let $M$ be a smooth manifold with a transitive action of a Lie group $G$. For any $g\in G$ consider the left translation $\tau(g)(p)=g\cdot p$ for all $p\in $. Let $X$ denote a vector field on $M$. Then one can define vector field $\tau_g(X)$ by the formula 
$$(\tau_{g}(X))_p=(\tau_g)_*(X_{g^{-1}\cdot p}).$$
An invariant affine connection on $M$ is an affine connection $\nabla$ which satisfies 
$$\tau_g(\nabla_XY)=\nabla_{\tau_g(X)}\tau_g(Y).$$
Recall that for every affine connection $\nabla$ and an arbitrary vector field $Z$ there is a $(1,1)$-tensor field on $M$ given by
$$L_Z^{\nabla}X=[Z,X]-\nabla_ZX.$$
Since $G$ acts transitively on $M$ and $\nabla$ is $G$-invariant, we get a commutative diagram
$$
\CD
T_oM @>{L^{\nabla}_{(Ad g^{-1}A)^+}}>> T_oM\\
@A{ (\tau_{g^{-1}})_*}AA  @A{(\tau_{g^{-1})_*}}AA\\
T_oM @>>{L^{\nabla}_{A^+}}> T_oM
\endCD
$$
Let $G/H$ be a homogeneus space. We will call it reductive with a reductive decomposition $\mathfrak{g}=\mathfrak{h}\oplus\mathfrak{m}$, if $Ad(H)(\mathfrak{m})\subset\mathfrak{m}$. Note that we fix the reductive decomposition, since it is not unique and the same manifold $M$ may have different presentations as a homogeneous space $G/H$.  For the given reductive decomposition one can identify $\mathfrak{m}$ with $T_oM$. In this case  the above diagram can be rewritten as 
$$
\CD
T_oM @>{L^{\nabla}_{(Ad g^{-1}A)^+}}>> T_oM\\
@A{\pi_*}AA @A{\pi_*}AA\\
\mathfrak{m} @>>{\alpha_{\nabla}(A,-)}> \mathfrak{m}
\endCD
$$
%where $\pi: G\rightarrow G/H$ is a natural projection. From this diagram we obtain
%$$\nabla_{A^+_p}B^+=[A^+,B^+]_p_*-((\tau_g)_*(\pi_*(\alpha_{\nabla}(Ad g^{-1}A,Ad g^{-1}B))).$$
This formula shows that $\nabla$ is recovered from the bilinear $Ad(H)$-equivariant map 
$$\alpha_{\nabla}:\mathfrak{m}\times\mathfrak{m}\rightarrow \mathfrak{m}.$$
For more details we refer to \cite{DGP}.
Finally we will formulate the above considerations as follows.
\begin{theorem}[Nomizu, \cite{N}] Let $G/H$ be a reductive homogeneous space with a fixed reductive decomposition $\mathfrak{g}=\mathfrak{h}\oplus\mathfrak{m}$. There is a bijective correspondence between the set of (left) $G$-invariant connections $\nabla$ on $G/H$ and the vector space of $Ad(H)$-equivariant  bilinear maps 
$$\alpha:\mathfrak{m}\times\mathfrak{m}\rightarrow\mathfrak{m}.$$
\end{theorem}
\noindent In the sequel we will call $\alpha$ the {\it Nomizu function}.


\section{Weyl curvature of left-invariant Weyl connections}
Let $(G,\langle -,-\rangle)$ be a Lie group endowed with an invariant Riemannian metric.
\begin{proposition}[\cite{KN}, Chapter X, Theorem 3.3]\label{prop:KN-LC} The Nomizu function $\alpha$ of the Levi-Civita connection on $(G,\langle,-,\rangle)$ is given by the formula
$$\alpha(X,Y)={1\over 2}[X,Y]+U(X,Y)$$
where $U(X,Y)$ is determined by
$$2\langle U(X,Y),Z\rangle=\langle X,[Z,Y]\rangle+\langle [Z,X],Y\rangle.$$
\end{proposition}
\begin{proposition}\label{prop:hat-k} Let $(G,\langle -,-\rangle)$ be a Lie group endowed with a left-invariant Weyl connection determined by $E\in\mathfrak{g}$. The Weyl curvature of this connection is given by the formula
 \begin{equation}\widehat{K}(X,Y) =$$
$$ K(X,Y)  + \langle {Y},{E}\rangle^2 + \langle {X},{E}\rangle^2 - \langle{E},{E}\rangle - \langle{Y},{\alpha(Y,E)}\rangle  - \langle{\alpha(X,E)},{X}\rangle,
\end{equation}
for any unit orthogonal $X,Y$.
\end{proposition}
\begin{proof} Use the following formula for the Weyl curvature \cite{W} (Proposition 3.2):
$$\widehat{R}_a(X,Y)Z=R(X,Y)Z+\langle Z,E\rangle(\langle Y,E\rangle X-\langle X,E\rangle Y)$$
$$+(\langle Z,Y\rangle \langle X,E\rangle-\langle Z,X\rangle\langle Y,E\rangle)E+E^2(\langle Z,X\rangle Y-\langle Z,Y\rangle X)$$
$$+\langle Z,\nabla_XE\rangle Y-\langle Z,\nabla_YE\rangle X+\langle Z,X\rangle\nabla_YE-\langle Z,Y\rangle\nabla_XE$$ 
  together with assumptions $Z=Y$, $\langle X,X\rangle=\langle Y,Y\rangle=1$, $\langle X,Y\rangle=0$ and the formula $(\nabla_UV)_o=\alpha(U,V)$.
\end{proof}

  

\section{Unimodular 3-dimensional solvable Lie groups with SNP Weyl connections}

The purpose of this section is to prove Theorem \ref{thm:main}.
\subsection{Left-invariant Riemannian metrics on completely solvable Lie groups}
\begin{definition}[\cite{A}]\label{thm:aleks} {\rm  A Lie algebra $\mathfrak{g}$ together with a fixed scalar product $\langle-,-\rangle$ is called a metric Lie algebra. Two metric Lie algebras $\mathfrak{g}$ and $\mathfrak{g}'$ are called isomorphic if there exists an isometry  $\varphi:\mathfrak{g}\rightarrow\mathfrak{g}'$ which  is simultaneously  an isomorphism of Lie algebras.} 
\end{definition}
As an example we consider solvable Lie algebras which have the following property: for any $X\in\mathfrak{g}$ the operator $\operatorname{ad}_X$ has only real eigenvalues.  
\begin{theorem}[\cite{A}]\label{thm:aleks} Isometric completely solvable metric Lie algebras are isomorphic.
\end{theorem}



\subsection{Classification of left-invariant Riemannian metris and sectional curvaures on 3-dimensional Lie groups}
We explain the classification of left-invariant metrics on Lie groups up to isomorphism. Let $G$ be a simply connected Lie group and $l_g: G\rightarrow G$ the left translation determined by $g$. A left invariant Riemannian metric on $G$ yields a scalar product on the Lie algebra $\mathfrak{g}$. Conversely, if $\langle -,-\rangle$ is a scalar product on $\mathfrak{g}$, the formula
$$\langle v,e\rangle_g=\langle (l_{g^{-1}})_*v,(l_{g^{-1}})_*w\rangle_e$$ 
defines a left-invariant Riemannian metric on $G$. The group $\operatorname{Aut}(\mathfrak{g})$ acts on the set $\mathcal{M}$ of all scalar products on $\mathfrak{g}$ by the formula
$$\varphi\cdot\langle v,w\rangle=\langle\varphi^{-1}v,\varphi^{-1}w\rangle.$$
Note that if one defines left-invariant Riemannian metrics on $G$ by the scalar products $\langle -,-\rangle$ and $\varphi\cdot\langle -,-\rangle$  lying in the same $\operatorname{Aut}(\mathfrak{g})$-orbit, then $(G,\langle -,-\rangle)$ and $(G,\varphi\cdot\langle v,w\rangle)$ will be isometric. Thus, a description of the space $\mathcal{M}/\operatorname{Aut}(\mathfrak{g})$ is a classification of left-invariant metrics on simply connected Lie groups {\it up to isomorphism}. In general not every isometry of $G$ is determined by an automorphism of $\mathfrak{g}$, thus this classification is finer then the classification up to isometry. However, it is  sufficient for our purposes. In our calculations we use this clasification in the form \cite{KHL}. Here is the sum up of this description. Let $[g]=(g_{ij})$ and $[g']=[g'_{ij}]$ be positive definite matrices representing the scalar products on $\mathfrak{g}$. We write $[g]\simeq [g']$ if 
$$[g']=[\varphi]^t[g][\varphi].$$
In this form, the classification of left-invariant Riemannian metrics on $G$ amounts to finding matrices representing a scalar product in each $\operatorname{Aut}(\mathfrak{g})$-orbit. Let us write down these matrices in the case of solvable unimodular 3-dimensional Lie groups.
It is well known that such groups are
\begin{enumerate}
\item abelian $\mathbb{R}^3$ (which is clear and not interesting from our point of view),
\item the Heisenberg group NIL whose Lie algebra is given by the Lie bracket formulas for the base $\{X,Y,Z\}$
$$[X,Y]=Z,[Z,X]=[Z,Y]=0.$$
\item the solvable group SOL, whose Lie algebra is given by
$$[X,Y]=0,[Z,X]=X,[Z,Y]=-Y$$
\item the universal cover  $\widetilde{E(2)}_0$ of the Lie group of proper motions of the euclidean plane  with the Lie algebra
$$[X,Y]=0,[Z,X]=-Y,[Z,Y]=X.$$

\end{enumerate}
The matrices representing classes $[g]$ are
\begin{enumerate}
\item the Heisenberg group NIL
$$
\begin{pmatrix}
\lambda & 0 & 0\\
0 & \lambda & 0\\
0 & 0 &1
\end{pmatrix}, \lambda>0
$$
\item the group SOL
$$
\begin{pmatrix}
1 & 0 & 0\\
0 & 1 & 0\\
0 & 0 & \nu
\end{pmatrix}, 
\begin{pmatrix}
1 & 1 & 0\\
1 & \mu & 0\\
0 & 0 & \nu
\end{pmatrix}, \mu>1, \nu>0
$$
\item the solvable Lie group $\tilde{E}_0(2)$
$$
\begin{pmatrix}
1 & 0  & 0\\
0 & \mu & 0\\
0 & 0 & \nu
\end{pmatrix}, 0<\mu \leq 1,\nu >0.
$$
 
\end{enumerate} 
In our approach we also need an explicit classification of $\operatorname{Aut}(\mathfrak{g})$ from \cite{KHL}.
\begin{enumerate}
\item For the Heisenberg algebra $\mathfrak{g}$ the group $\operatorname{Aut}(\mathfrak{g})$ is isomorphic to a matrix group
consiting of matrices
$$
\begin{pmatrix}
a & c & 0\\
b & d & 0\\
* & * & ad-bc\\
\end{pmatrix}, a,b,c,d\in\mathbb{R},ad-bc\not=0.
$$
\item For the Lie algebra SOL the group $\operatorname{Aut}(\mathfrak{g})$ is isomorphic to the group consisting of matrices
$$
S=\begin{pmatrix}
x_1 & 0 & x_3\\
0 &  y_2 & y_3\\
0 & 0 & 1\\
\end{pmatrix}, x_1,x_3,y_2,y_3\in\mathbb{R},\,\text{or}\, 
S\cdot
\begin{pmatrix}
0 & 1 & 0\\
1 & 0 & 0\\
0 & 0 & -1\\
\end{pmatrix}
$$
\item For the unimodular solvable Lie algebra $\mathfrak{g}=\mathbb{R}^2\rtimes\mathfrak{so}(2)$ the automorphism group $\operatorname{Aut}(\mathfrak{g})$ consists of matrices
$$
\begin{pmatrix}
a & b & \gamma\\
-b & a & \delta\\
0 & 0 & 1\\
\end{pmatrix}, (a,b)\not=(0,0), a,b, \gamma,\delta\in\mathbb{R}.
$$
\end{enumerate}
  
\subsection{Calculation of Ricci curvatures and sectional curvatures in the 3-dimensional case}

 By \cite{M}, Theorem 4.3,  on every 3-dimensional unimodular Lie group, an orthonormal basis $\{Y_1,Y_2,Y_3\}$ with \begin{align*}
    \begin{matrix}
        [Y_2,Y_3] = a_1Y_1, & [Y_3,Y_1] = a_2Y_2,& [Y_1,Y_2] = a_3Y_3
    \end{matrix}
\end{align*}
can be found. We will call such basis a \textit{Milnor basis}. In such basis, the Ricci tensor is diagonal: indeed, in unimodular situation,
\begin{align*}
    \text{Ric}(X,Y) = &-\dfrac{1}{2}\sum_{k}\langle [X,X_k],[Y,X_k] \rangle \\
    &+ \dfrac{1}{4}\sum_{k,l}\langle[X_k,X_l],X\rangle\langle[X_k,X_l],Y\rangle - \dfrac{1}{2}B(X,Y).
\end{align*}
Here $B(X,Y)$ is the Killing form. Let's conduct the calculations for $\text{Ric}(Y_1,Y_2)$. The first term vanishes: for $k=1,2$ it follows from anti-symmetry of the Lie bracket, and for $k=3$ we obtain \begin{align*}
    \langle [Y_1,Y_3],[Y_2,Y_3]\rangle = \langle -a_2Y_2,a_1Y_1\rangle = 0.
\end{align*} For the second term, notice that each bracket $[Y_k,Y_l]$ is a multiple of exactly one basis vector. In particular, at most one of the factors $\langle[Y_k,Y_l],Y_1\rangle$, $\langle[Y_k,Y_l],Y_2\rangle$ is nonzero. \\~\\
Similarly, we show that $B(Y_k,Y_l) = 0$ for $k \neq l$. Consider $B(Y_1,Y_2)$. In order to compute it, we calculate $\text{ad}_{Y_2}Y_1 = -a_3Y_3$, so  $\text{ad}_{Y_1}(-a_3Y_3) = a_2a_3Y_2$. Therefore $B(Y_1,Y_2) = \text{tr}(\text{ad}_Y \circ \text{ad}_Y) = 0$.\\~\\
The other cases follow by taking any cyclic permutation of the basic vector fields. Now we just need to compute the diagonal entries $\text{Ric}(Y_i,Y_i)$; again, we can just compute $\text{Ric}(Y_1,Y_1)$ as the other cases will follow by taking a cyclic permutation. \\~\\
We will use the same formula as before, which in this case simplifies to $$\text{Ric}(Y_1,Y_1) = -\dfrac{1}{2}\sum_{k}||[Y_1,Y_k]||^2 + \dfrac{1}{4}\sum_{k,l}\langle[Y_k,Y_l],Y_1\rangle^2 - \dfrac{1}{2}B(Y_1,Y_1).$$
For the first term, the only contributing factors are $[Y_1,Y_2] = a_3Y_3$ and $[Y_1,Y_3] = a_2Y_2$. It therefore evaluates to $-\dfrac{1}{2}\left(a_2^2 + a_3^2\right)$. \\~\\ For the second term, only the brackets $[Y_2,Y_3] = a_1Y_1$ and $[Y_3,Y_2] = -a_1Y_1$ are going to contribute; each one contributes with $a_1^2$. Therefore, the second term evaluates to $\dfrac{1}{4}\left(a_1^2+a_1^2\right) = \dfrac{1}{2}a_1^2$. \\~\\
Finally, $B(Y_1,Y_1) = -2a_2a_3$, so the third term evaluates to $\dfrac{1}{2}a_2a_3$. Combining them together, we obtain that \begin{align*}
    \text{Ric}(Y_1,Y_1) &= -\dfrac{1}{2}(a_2^2 + a_3^2) + \dfrac{1}{2}a_1^2 + \dfrac{1}{2}a_2a_3 \\ 
    &= \dfrac{1}{2}(a_1^2 - a_2^2 + 2a_2a_3 - a_3^2) = \dfrac{1}{2}(a_1^2 - (a_2-a_3)^2).
\end{align*} 
By cyclic permutations $(i,j,k)$, we obtain the formula for the principal Ricci curvatures
$$\rho_i = \text{Ric}(Y_i,Y_i) = \dfrac{1}{2}(a_i^2 - (a_j - a_k)^2).$$
Denoting the principal Ricci curvatures as $\rho_1,\rho_2,\rho_3$, one may compute the sectional curvatures $K(Y_i,Y_j)$ as \begin{align*}
    K(Y_i,Y_j) = \dfrac{\rho_i + \rho_j - \rho_k}{2}.
\end{align*}
Denote $K(Y_i,Y_j)$ as $K_{ij}$, so that \begin{align*}
    \rho_i &= K_{ij}+ K_{ik},\\
    \rho_j &= K_{ij} + K_{jk}, \\
    \rho_k &= K_{ik} + K_{jk}.
\end{align*}
Thus $\rho_i + \rho_j - \rho_k = 2K_{ij}$ and the formula follows immediately.

\subsection{Invariant Riemannian metric (I) for SOL}
\subsubsection{Sectional curvatures}
Let $\{X,Y,Z\}$ be the basis of the Lie algebra of SOL as in the classification found in \cite{KHL}, so \begin{align*}
    [X, Y] = 0, \qquad [Z,X] = X, \qquad [Z, Y] = -Y.
\end{align*}
By taking $$Y_1 = \dfrac{1}{\sqrt{2}}(X+Y),\qquad Y_2 = -\dfrac{1}{\sqrt{2}}(X-Y), \qquad Y_3 = \dfrac{1}{\sqrt{\nu}}Z, $$
we obtain an orthonormal Milnor base with $$[Y_2,Y_3] = \dfrac{1}{\sqrt{\nu}} Y_1,\; [Y_3,Y_1] = -\dfrac{1}{\sqrt{\nu}}Y_2,\; [Y_1,Y_2] = 0.$$
Therefore, the principal Ricci curvatures are \begin{align*}
    \rho_1 &= \dfrac{1}{2}\left(\dfrac{1}{\nu}-\dfrac{1}{\nu}\right) = 0, \\
    \rho_2 &= \dfrac{1}{2}\left(\dfrac{1}{\nu} - \dfrac{1}{\nu}\right) = 0, \\
    \rho_3 &= \dfrac{1}{2} \left(\dfrac{1}{\sqrt{\nu} } + \dfrac{1}{\sqrt{\nu}}\right)^2 = -\dfrac{2}{\nu}.
\end{align*}
Hence, the sectional curvatures $K(Y_i,Y_j)$ are given by \begin{align*}
    K(Y_1,Y_2) &=  \dfrac{1}{\nu}, \\
    K(Y_2,Y_3) &= -\dfrac{1}{\nu}, \\
    K(Y_3,Y_1) &= -\dfrac{1}{\nu}.
\end{align*}

\subsubsection{Weyl sectional curvature}

We  compute the terms $U(Y_i,Y_j)$ using Proposition \ref{prop:KN-LC}. For $U(Y_1,Y_2)$, the components along $Y_1$ and $Y_2$ vanish. The only possibly nonzero component is
\begin{align*}
    2\left\langle U(Y_1,Y_2),Y_3 \right\rangle
    &=\left\langle Y_1,[Y_3,Y_2] \right\rangle+\left\langle [Y_3,Y_1],Y_2 \right\rangle \\
    &=\left\langle Y_1,-\dfrac{1}{\nu}Y_1 \right\rangle+\left\langle -\dfrac{1}{\nu}Y_2,Y_2 \right\rangle
    =-\dfrac{2}{\nu}.
\end{align*}
Hence
\begin{align*}
    U(Y_1,Y_2)=-\dfrac{1}{\nu}Y_3.
\end{align*}
Similarly,
\begin{align*}
    2\left\langle U(Y_1,Y_3),Y_2 \right\rangle
    &=\left\langle Y_1,[Y_2,Y_3] \right\rangle+\left\langle [Y_2,Y_1],Y_3 \right\rangle=\dfrac{1}{\nu}, \\
    2\left\langle U(Y_2,Y_3),Y_1 \right\rangle
    &=\left\langle Y_2,[Y_1,Y_3] \right\rangle+\left\langle [Y_1,Y_2],Y_3 \right\rangle=\dfrac{1}{\nu}.
\end{align*}
Therefore
\begin{align*}
    U(Y_1,Y_2)=-\dfrac{1}{\nu}Y_3,
    \qquad
    U(Y_1,Y_3)=\frac{1}{2\nu}Y_2,
    \qquad
    U(Y_2,Y_3)=\frac{1}{2\nu}Y_1.
\end{align*}
Moreover, $U(Y_i,Y_i)=0$ for $i=1,2,3$. \\~\\
Therefore the values of the Nomizu function $\alpha$ are given by
\begin{align*}
\begin{matrix*}[l]
    \alpha(Y_1,Y_2)=-\dfrac{1}{\nu}Y_3, &&     \alpha(Y_1,Y_3)=\dfrac{1}{\nu}Y_2, &&  \alpha(Y_2,Y_3)=\dfrac{1}{\nu}Y_1, \\[1em]
    \alpha(Y_2,Y_1)=-\dfrac{1}{\nu}Y_3, && 
    \alpha(Y_3,Y_1)=0, &&
    \alpha(Y_3,Y_2)=0.
\end{matrix*}
\end{align*}
Also
\begin{align*}
    \alpha(Y_i,Y_i)=0,\qquad i=1,2,3.
\end{align*}

\noindent Let $E$ be the vector field defining the Weyl connection, given in the basis $Y_1,Y_2,Y_3$ by
\begin{align*}
    E=e_1Y_1+e_2Y_2+e_3Y_3.
\end{align*}
As the basis $\{Y_1,Y_2,Y_3\}$ is orthonormal, we get that
\begin{align*}
    \left\langle Y_1,E \right\rangle=e_1,
    \qquad
    \left\langle Y_2,E \right\rangle=e_2,
    \qquad
    \left\langle Y_3,E \right\rangle=e_3,
    \qquad
    \left\langle E,E \right\rangle=e_1^2+e_2^2+e_3^2.
\end{align*}
and find
\begin{align*}
    \nonumber \alpha(Y_1,E)
    &=e_1\alpha(Y_1,Y_1)+e_2\alpha(Y_1,Y_2)+e_3\alpha(Y_1,Y_3) \\
    &=-\dfrac{1}{\nu}e_2Y_3+\dfrac{1}{\nu}e_3Y_2,
\end{align*}
\begin{align*}
    \nonumber \alpha(Y_2,E)
    &=e_1\alpha(Y_2,Y_1)+e_2\alpha(Y_2,Y_2)+e_3\alpha(Y_2,Y_3) \\
    &=-\dfrac{1}{\nu}e_1Y_3+\dfrac{1}{\nu}e_3Y_1,
\end{align*}
\begin{align*}
    \nonumber \alpha(Y_3,E)
    &=e_1\alpha(Y_3,Y_1)+e_2\alpha(Y_3,Y_2)+e_3\alpha(Y_3,Y_3) \\
    &=0.
\end{align*}
In particular,
\begin{align*}
    \left\langle \alpha(Y_i,E),Y_i \right\rangle=0,
    \qquad i=1,2,3.
\end{align*}
We now substitute these values into the formula for the Weyl curvature given by Proposition \ref{prop:hat-k}. For the plane spanned by $Y_1,Y_2$, we have
\begin{align*}
    \nonumber \widehat{K}(Y_1,Y_2)
    &=K(Y_1,Y_2)+e_2^2+e_1^2-(e_1^2+e_2^2+e_3^2) \\
    &\qquad -\left\langle Y_2,\alpha(Y_2,E) \right\rangle
    -\left\langle \alpha(Y_1,E),Y_1 \right\rangle \\
    &=\dfrac{1}{\nu}-e_3^2.
\end{align*}
Analogously
\begin{align*}
    \widehat{K}(Y_2,Y_3)
    &=-\frac{1}{\nu}-e_1^2, \\
    \widehat{K}(Y_3,Y_1) &= \frac{1}{\nu}-e_2^2.
\end{align*}

\subsection{Invariant Riemannian metric (II) for SOL}

\subsubsection{Sectional curvatures}
Note that the calculation in \cite{KHL} contains a small inaccuracy, and for the sake of correctness, we recalculate the sectional curvatures, although it does not influence the final results  neither in \cite{KHL} nor in this article.
 The basis $X,Y,Z$ of the Lie algebra again satisfies
\begin{align*}
    [X,Y] = 0, \; [Z,X] = X,\; [Z,Y] = -Y.
\end{align*}
We take the orthonormal Milnor base
$$
Y_1 =
\dfrac{\sqrt{\mu}X+Y}{\sqrt{2}\sqrt{\mu+\sqrt{\mu}}},
\qquad
Y_2 =
\dfrac{-\sqrt{\mu}X+Y}{\sqrt{2}\sqrt{\mu-\sqrt{\mu}}},
\qquad
Y_3 =
\dfrac{1}{\sqrt{\nu}}Z
$$
with the following easily calculated structure constants:
$$
[Y_2,Y_3]
=
\dfrac{\sqrt{\mu+\sqrt{\mu}}}{\sqrt{\nu}\sqrt{\mu-\sqrt{\mu}}}Y_1,
\qquad
[Y_3,Y_1]
=
-\dfrac{\sqrt{\mu-\sqrt{\mu}}}{\sqrt{\nu}\sqrt{\mu+\sqrt{\mu}}}Y_2,
\qquad
[Y_1,Y_2]=0.
$$
For clarity, denote
$$
a_1=
\dfrac{\sqrt{\mu+\sqrt{\mu}}}{\sqrt{\nu}\sqrt{\mu-\sqrt{\mu}}},
\qquad
a_2=
-\dfrac{\sqrt{\mu-\sqrt{\mu}}}{\sqrt{\nu}\sqrt{\mu+\sqrt{\mu}}},
\qquad
a_3=0.
$$
 We compute the principal Ricci curvatures:
\begin{align*}
    \rho_1
    &=
    \dfrac{1}{2}\left(a_1^2-(a_2-a_3)^2\right) =
    \dfrac{1}{2}\left(
    \dfrac{\mu+\sqrt{\mu}}{\nu(\mu-\sqrt{\mu})}
    -
    \dfrac{\mu-\sqrt{\mu}}{\nu(\mu+\sqrt{\mu})}
    \right) \\
    &=
    \dfrac{2\sqrt{\mu}}{\nu(\mu-1)}, \\[1em]
    \rho_2
    &=
    \dfrac{1}{2}\left(a_2^2-(a_3-a_1)^2\right) =
    \dfrac{1}{2}\left(
    \dfrac{\mu-\sqrt{\mu}}{\nu(\mu+\sqrt{\mu})}
    -     \dfrac{\mu+\sqrt{\mu}}{\nu(\mu-\sqrt{\mu})}
    \right) \\
    &=
    -\dfrac{2\sqrt{\mu}}{\nu(\mu-1)}, \\[1em]
    \rho_3
    &=
    \dfrac{1}{2}\left(a_3^2-(a_1-a_2)^2\right) =
    -\dfrac{1}{2}
    \left(
    \dfrac{\sqrt{\mu+\sqrt{\mu}}}{\sqrt{\nu}\sqrt{\mu-\sqrt{\mu}}}
    +
    \dfrac{\sqrt{\mu-\sqrt{\mu}}}{\sqrt{\nu}\sqrt{\mu+\sqrt{\mu}}}
    \right)^2 \\
    &=
    -\dfrac{2\mu}{\nu(\mu-1)}.
\end{align*}
Therefore
\begin{align*}
    K(Y_1,Y_2)
    &=
    \dfrac{\rho_1+\rho_2-\rho_3}{2} =
    \dfrac{1}{2}\left(
    \dfrac{2\sqrt{\mu}}{\nu(\mu-1)}
    -
    \dfrac{2\sqrt{\mu}}{\nu(\mu-1)}
    +
    \dfrac{2\mu}{\nu(\mu-1)}\right) \\
    &=
    \dfrac{\mu}{\nu(\mu-1)}, \\[1em]
    K(Y_2,Y_3)
    &=
    \dfrac{\rho_2+\rho_3-\rho_1}{2} =
    \dfrac{1}{2}\left(
    -\dfrac{2\sqrt{\mu}}{\nu(\mu-1)}
    -
    \dfrac{2\mu}{\nu(\mu-1)}
    -
    \dfrac{2\sqrt{\mu}}{\nu(\mu-1)}\right) \\
    &=
    -\dfrac{\mu+2\sqrt{\mu}}{\nu(\mu-1)}, \\[1em]
    K(Y_3,Y_1)
    &=
    \dfrac{\rho_3+\rho_1-\rho_2}{2} =
    \dfrac{1}{2}\left(
    -\dfrac{2\mu}{\nu(\mu-1)}
    +
    \dfrac{2\sqrt{\mu}}{\nu(\mu-1)}
    +
    \dfrac{2\sqrt{\mu}}{\nu(\mu-1)}\right) \\
    &=
    \dfrac{2\sqrt{\mu}-\mu}{\nu(\mu-1)}.
\end{align*}
Now we calculate the Weyl curvatures by the same method as in case (I).
Starting with the bilinear function $U(-,-)$, the only possibly nonzero term in this case it
\begin{align*}
    2\left\langle U(Y_1,Y_2),Y_3 \right\rangle
    &=\left\langle Y_1,[Y_3,Y_2] \right\rangle+\left\langle [Y_3,Y_1],Y_2 \right\rangle \\
    &=\left\langle Y_1,-a_1Y_1 \right\rangle+\left\langle a_2Y_2,Y_2 \right\rangle
    =a_2-a_1.
\end{align*}
Hence
\begin{align*}
    U(Y_1,Y_2)=\frac{a_2-a_1}{2}Y_3.
\end{align*}
Similarly,
\begin{align*}
    2\left\langle U(Y_1,Y_3),Y_2 \right\rangle
    &=\left\langle Y_1,[Y_2,Y_3] \right\rangle+\left\langle [Y_2,Y_1],Y_3 \right\rangle=a_1, \\
    2\left\langle U(Y_2,Y_3),Y_1 \right\rangle
    &=\left\langle Y_2,[Y_1,Y_3] \right\rangle+\left\langle [Y_1,Y_2],Y_3 \right\rangle=-a_2.
\end{align*}
Therefore
\begin{align*}
    U(Y_1,Y_2)=\frac{a_2-a_1}{2}Y_3,
    \qquad
    U(Y_1,Y_3)=\frac{a_1}{2}Y_2,
    \qquad
    U(Y_2,Y_3)=-\frac{a_2}{2}Y_1.
\end{align*}
Moreover, $U(Y_i,Y_i)=0$ for $i=1,2,3$. \\~\\ Therefore values of the Nomizu function $\alpha$ are given by
\begin{align*}
    \begin{matrix*}[l]
      \alpha(Y_1,Y_2)=\dfrac{a_2-a_1}{2}Y_3, && \alpha(Y_1,Y_3)=\dfrac{a_1-a_2}{2}Y_2, && \alpha(Y_2,Y_3)=\dfrac{a_1-a_2}{2}Y_1, 
       \\[1em]
       \alpha(Y_2,Y_1)=\dfrac{a_2-a_1}{2}Y_3,&&
      \alpha(Y_3,Y_1)= \dfrac{a_1+a_2}{2}Y_2, && \alpha(Y_3,Y_2)=-\dfrac{a_1+a_2}{2}Y_1
    \end{matrix*}
\end{align*}
and, as previously,
\begin{align*}
    \alpha(Y_i,Y_i)=0,\qquad i=1,2,3.
\end{align*}
Taking
\begin{align*}
    E=e_1Y_1+e_2Y_2+e_3Y_3,
\end{align*}
 we find
\begin{align*}
      \alpha(Y_1,E)
    &=e_1\alpha(Y_1,Y_1)+ e_2\alpha(Y_1,Y_2)+ e_3\alpha(Y_1,Y_3) \\
    &=\frac{a_2-a_1}{2}e_2Y_3+\frac{a_1-a_2}{2}e_3Y_2, \\[1em]
    \alpha(Y_2,E)
    &=e_1\alpha(Y_2,Y_1)+ e_2\alpha(Y_2,Y_2)+ e_3\alpha(Y_2,Y_3) \\
    &=\frac{a_2-a_1}{2}e_1Y_3+\frac{a_1-a_2}{2}e_3Y_1, \\[1em]
    \alpha(Y_3,E)
    &=e_1\alpha(Y_3,Y_1)+ e_2\alpha(Y_3,Y_2)+ e_3\alpha(Y_3,Y_3) \\
    &=\frac{a_1+a_2}{2}e_1Y_2-\frac{a_1+a_2}{2}e_2Y_1
\end{align*}
In particular,
\begin{align*}
    \left\langle \alpha(Y_i,E),Y_i \right\rangle=0,
    \qquad i=1,2,3.
\end{align*}
Finally, we have
\begin{align*}
    \nonumber \widehat{K}(Y_1,Y_2)
    &=K(Y_1,Y_2)+e_2^2+e_1^2-(e_1^2+e_2^2+e_3^2) \\
    &\qquad \qquad \qquad -\left\langle Y_2,\alpha(Y_2,E) \right\rangle
    -\left\langle \alpha(Y_1,E),Y_1 \right\rangle \\
    &=\frac{\mu}{\nu(\mu-1)}-e_3^2.
\end{align*}
Analogously
\begin{align*}
    \widehat{K}(Y_2,Y_3) =-\frac{\mu+2\sqrt{\mu}}{\nu(\mu-1)}-e_1^2, \\[1em]
    \widehat{K}(Y_3,Y_1) =\frac{2\sqrt{\mu}-\mu}{\nu(\mu-1)}-e_2^2.
\end{align*}
\subsection{Heisenberg algebra}
\noindent The standard basis $\{X,Y,Z\}$ of the Lie algebra of the Heisenberg group $\mathrm{Nil}$ satisfies
\begin{align*}
\noindent [X,Y]=Z,
\qquad
[Z,X]=[Z,Y]=0.
\end{align*}

\noindent We consider the orthonormal Milnor base
\begin{align*}
\noindent
Y_1=\frac{1}{\sqrt{\lambda}}X,
\qquad
Y_2=\frac{1}{\sqrt{\lambda}}Y,
\qquad
Y_3=Z.
\end{align*}

\noindent with the structure constants
\begin{align*}
\noindent
[Y_1,Y_2]=\frac{1}{\lambda}Y_3,
\qquad
[Y_2,Y_3]=0,
\qquad
[Y_3,Y_1]=0.
\end{align*}

\noindent The sectional curvatures are
\begin{align*}
K(Y_1,Y_2)=-\frac{3}{4\lambda^2},
\qquad
K(Y_2,Y_3)=\frac{1}{4\lambda^2},
\qquad
K(Y_3,Y_1)=\frac{1}{4\lambda^2}.
\end{align*}

\noindent The values of the terms $U(Y_1,Y_j),\; i=1,2,3$ are
\begin{align*}
U(Y_1,Y_2)=0,
\qquad
U(Y_1,Y_3)=-\frac{1}{2\lambda}Y_2,
\qquad
U(Y_2,Y_3)=\frac{1}{2\lambda}Y_1.
\end{align*}

\noindent Also,
\begin{align*}
U(Y_i,Y_i)=0
\qquad
\text{for } i=1,2,3.
\end{align*}

\noindent Thus the relevant values of the Nomizu function $\alpha$ are
\begin{align*}
\begin{matrix*}[l]
    \alpha(Y_1,Y_2)=\dfrac{1}{2\lambda}Y_3, &&
\alpha(Y_1,Y_3)=-\dfrac{1}{2\lambda}Y_2, && \alpha(Y_2,Y_3)=\dfrac{1}{2\lambda}Y_1, \\[1em]
\alpha(Y_2,Y_1)=-\dfrac{1}{2\lambda}Y_3, &&
\alpha(Y_3,Y_1)=-\dfrac{1}{2\lambda}Y_2, &&
\alpha(Y_3,Y_2)=\dfrac{1}{2\lambda}Y_1.
\end{matrix*}
\end{align*}

\noindent For any left-invariant vector field $E = e_1Y_1 + e_2Y_2 + e_3Y_3,$ the values of $\alpha(Y_i,E)$ are given by
\begin{align*}
\alpha(Y_1,E) &= \dfrac{e_2}{2\lambda}Y_2 - \dfrac{e_3}{2\lambda}Y_3, \\
\alpha(Y_2,E) &= \dfrac{e_3}{2\lambda}Y_1 -\dfrac{e_1}{2\lambda}Y_3 , \\
\alpha(Y_3,E) &=  \dfrac{e_2}{2\lambda}Y_1-\dfrac{e_1}{2\lambda}Y_2.
\end{align*}

\noindent In particular,
\begin{align*}
\noindent
\langle \alpha(Y_i,E),Y_i\rangle=0
\qquad
\text{for } i=1,2,3.
\end{align*}

\noindent Hence
\begin{align*}
\widehat{K}(Y_1,Y_2) &=
-\frac{3}{4\lambda^2}-e_3^2. \\
\widehat{K}(Y_2,Y_3) &=
\frac{1}{4\lambda^2} - e_1^2, \\
\widehat{K}(Y_3,Y_1) &=
\frac{1}{4\lambda^2} - e_2^2.
\end{align*}

\subsection{$\mathbb{R}^2 \rtimes \mathfrak{so}(2)$}

\noindent The Lie algebra
$\mathbb{R}^2\rtimes\mathfrak{so}(2)$ of the universal cover $\mathbb{R}^2\rtimes SO(2)$ of  admits the basis $\{X,Y,Z\}$ satisfying
\begin{align*}
[X,Y]=0,\qquad [Z,X]=-Y,\qquad [Z,Y]=X.
\end{align*}
Consider the orthonormal Milnor base given by
\begin{align*}
Y_1=X,
\qquad
Y_2=\frac{1}{\sqrt{\mu}}Y,
\qquad
Y_3=\frac{1}{\sqrt{\nu}}Z
\end{align*}
with the structure constants
\begin{align*}
[Y_2,Y_3]
=
-\frac{1}{\sqrt{\mu\nu}}Y_1,
\qquad
[Y_3,Y_1]
=
-\frac{\sqrt{\mu}}{\sqrt{\nu}}Y_2,
\qquad
[Y_1,Y_2]=0.
\end{align*}
\noindent The values of the terms $U(Y_1,Y_j),\; i=1,2,3$ are
\begin{align*}
U(Y_1,Y_2)= \frac{1-\mu}{2\sqrt{\mu\nu}}Y_3,
\qquad
U(Y_1,Y_3)=-\frac{1}{2\sqrt{\mu\nu}}Y_2, \qquad 
U(Y_2,Y_3)= \frac{\mu}{2\sqrt{\mu\nu}}Y_1.
\end{align*}
with \begin{align*}
    U(Y_i,Y_i) = 0 \qquad \text{for}\; i=1,2,3.
\end{align*}
We then find the values $\alpha(Y_i,Y_j)$ of the Nomizu function:
\begin{align*}
\begin{matrix*}[l]
\alpha(Y_1,Y_2)=\dfrac{1-\mu}{2\sqrt{\mu\nu}}Y_3, &&
\alpha(Y_1,Y_3)=-\dfrac{1-\mu}{2\sqrt{\mu\nu}}Y_2, &&
\alpha(Y_2,Y_3)=-\dfrac{1-\mu}{2\sqrt{\mu\nu}}Y_1. \\[1em]
\alpha(Y_2,Y_1)=\dfrac{1-\mu}{2\sqrt{\mu\nu}}Y_3, &&
\alpha(Y_3,Y_1)=-\dfrac{1+\mu}{2\sqrt{\mu\nu}}Y_2, &&
\alpha(Y_3,Y_2)=\dfrac{1+\mu}{2\sqrt{\mu\nu}}Y_1.
\end{matrix*}
\end{align*}
Taking the left-invariant vector field $E = e_1Y_1 + e_2Y_2 + e_3Y_3$, we obtain
\begin{align*}
\alpha(Y_1,E) &= 
\frac{1-\mu}{2\sqrt{\mu\nu}}
\left(e_2Y_3-e_3Y_2\right), \\
\alpha(Y_2,E) &=
\frac{1-\mu}{2\sqrt{\mu\nu}}
\left(e_1Y_3-e_3Y_1\right), \\
\alpha(Y_3,E) &=
\frac{1+\mu}{2\sqrt{\mu\nu}}
\left(-e_1Y_2+e_2Y_1\right).
\end{align*}

\noindent Therefore
\begin{align*}
\langle \alpha(Y_i,E),Y_i\rangle=0
\qquad
\text{for } i=1,2,3.
\end{align*}
Finally, we obtain
\begin{align*}
\widehat{K}(Y_1,Y_2) &= 
\frac{(1-\mu)^2}{4\mu\nu} - e_3^2, \\
\widehat{K}(Y_2,Y_3) &= 
-\frac{(1-\mu)(3+\mu)}{4\mu\nu} -e_1^2, \\
\widehat{K}(Y_3,Y_1) &=
\frac{(1-\mu)(1+3\mu)}{4\mu\nu} - e_2^2.
\end{align*}

\subsection{Completion of proof}
We have proved that for the Milnor bases $\{Y_1,Y_2,Y_3\}$ the Weyl curvatures $\widehat{K}(Y_i,Y_j)\leq 0$, after the scaling of $E$, for all 3-dimensional unimodular solvable Lie algebras.  Now we restrict our considerations to the completely solvable case, that is to the groups NILand SOL. We  observe the following. By Theorem \ref{thm:aleks}  in each of these two cases the classification of left-invariant Riemannian metrics up to automorphism is equivalent to the classification  up to isometry. Therefore, if we change the base to another orthonormal base $\{Y_1',Y_2',Y_3'\}$ we can find an automorphism $\varphi\in\operatorname{Aut}(\mathfrak{g})$  in the isometry class which sends $Y_i$ to $Y_i'$.
Now we neeed to consider each of these cases separately.

\subsubsection{Group SOL}
\vskip6pt
{\bf Metric (II) on Sol}

Consider the left-invariant metric of type (II) on Sol. We first restrict
attention to automorphisms in the connected component of the identity. By the
description of \(\operatorname{Aut}(\mathfrak g)\), every such automorphism is
of the form
\[
\varphi(Y_1)=x_1Y_1,\qquad
\varphi(Y_2)=y_2Y_2,\qquad
\varphi(Y_3)=x_3Y_1+y_3Y_2+Y_3,
\]
where \(x_1,y_2\neq 0\). Since we are in the connected component of the
identity, we may assume \(x_1,y_2>0\).

Let $E=e_3Y_3.$ Then
\[ \varphi(E)=e_3\varphi(Y_3) =e_3x_3Y_1+e_3y_3Y_2+e_3Y_3. \]
Thus, with respect to the original Milnor base, the components of
\(\varphi(E)\) are
\[e_1'=e_3x_3,\qquad e_2'=e_3y_3,\qquad e_3'=e_3. \]

Recall that for the Milnor base  on SOL we have
\[
[Y_2,Y_3]=a_1Y_1,\qquad
[Y_3,Y_1]=a_2Y_2,\qquad
[Y_1,Y_2]=0.
\]

As \(e_1=e_2=0\), the Weyl sectional curvatures can be easily computed:
\[
\widehat K(\varphi(Y_1),\varphi(Y_2))=K_{12}-e_3^2,
\]
\[
\widehat K(\varphi(Y_2),\varphi(Y_3))=\frac{x_3^2(K_{12}-e_3^2)+K_{23}}{1+x_3^2},
\]
\[
\widehat K(\varphi(Y_3),\varphi(Y_1))=\frac{y_3^2(K_{12}-e_3^2)+K_{31}}{1+y_3^2}.
\]
Thus these curvatures are non-positive for every \(x_3,y_3\) if and only if
\[
K_{12}-e_3^2\leq 0,
\qquad
K_{23}\leq 0,
\qquad
K_{31}\leq 0.
\]
The second condition is automatically satisfied: substituting the values of $K_{12}$ and $K_{23}$ yields
\[
-\frac{(a_1-a_2)(3a_1+a_2)}{4}\leq 0.
\]
Since \(a_1-a_2>0\), this is equivalent to $3a_1+a_2\geq 0$,
which always holds. The first condition says
\[
|e_3|\geq \frac{a_1-a_2}{2}.
\]
and cannot be simplified further. The third condition is
\[
\frac{(a_1-a_2)(a_1+3a_2)}{4}\leq 0.
\]
Since \(a_1-a_2>0\), this is equivalent to $a_1+3a_2\leq 0,$ or equivalently to $\dfrac{-a_2}{a_1} \geq \dfrac{1}{3}.$ Substituting the values of $a_1$ and $a_2$ in terms of $\mu,\;\nu$ we obtain
\[
-\frac{a_2}{a_1}
=
\frac{\mu-\sqrt{\mu}}{\mu+\sqrt{\mu}}
=
\frac{\sqrt{\mu}-1}{\sqrt{\mu}+1}.
\]
Therefore we obtain
\[
\frac{\sqrt{\mu}-1}{\sqrt{\mu}+1}\geq \frac13.
\]
Solving this inequality in terms of $\mu$, we obtain that $\mu \geq 4.$ \\~\\
Thus the result is as follows: \begin{itemize}
    \item $\mu < 4$: no SNP connections,
    \item $\mu = 4$: two choices $E = \pm \dfrac{a_1-a_2}{2}Y_3$,
    \item $\mu > 4$: the two rays $|e_3| \geq \dfrac{a_1-a_2}{2}.$
\end{itemize}

Also, if one takes a matrix from the other component of Aut($\mathfrak{g}$), the calculations will be the same, only with interchanging vectors $Y_1$ and $Y_2$ and changing $E$ to $-E$. But the non-positivity won't change anyway, as changing $E$ to $-E$ won't change the terms $e_i^2$.

\subsection*{Metric (I) on Sol}

Now consider the left-invariant metric of type (I) on SOL, given by
\[
\begin{pmatrix}
1 & 0 & 0\\
0 & 1 & 0\\
0 & 0 & \nu
\end{pmatrix},
\qquad \nu>0.
\]
For the corresponding Milnor base we have
\[
[Y_2,Y_3]=a_1Y_1,\qquad
[Y_3,Y_1]=a_2Y_2,\qquad
[Y_1,Y_2]=0,
\]
where
\[
a_1=\frac{1}{\sqrt{\nu}},
\qquad
a_2=-\frac{1}{\sqrt{\nu}}.
\]
As before, we restrict attention to automorphisms in the connected component
of the identity. Thus
\[
\varphi(Y_1)=x_1Y_1,\qquad
\varphi(Y_2)=y_2Y_2,\qquad
\varphi(Y_3)=x_3Y_1+y_3Y_2+Y_3,
\]
where \(x_1,y_2>0\). As before, consider $E = e_3Y_3$. \\~\\

For the metric of type (I), the Riemannian sectional curvatures are
\[
K_{12}=\frac{1}{\nu}, \qquad 
K_{23} = 
-\frac{1}{\nu}, \qquad 
K_{31} =-\frac{1}{\nu}.
\]
Hence
\[
\widehat K(\varphi(Y_1),\varphi(Y_2))
=
K_{12}-e_3^2
=
\frac{1}{\nu}-e_3^2.
\]
Moreover,
\[
\widehat K(\varphi(Y_2),\varphi(Y_3))
=
\frac{x_3^2(K_{12}-e_3^2)+K_{23}}{1+x_3^2}
=
\frac{x_3^2\left(\frac{1}{\nu}-e_3^2\right)-\frac{1}{\nu}}
{1+x_3^2},
\]
and
\[
\widehat K(\varphi(Y_3),\varphi(Y_1))
=
\frac{y_3^2(K_{12}-e_3^2)+K_{31}}{1+y_3^2}
=
\frac{y_3^2\left(\frac{1}{\nu}-e_3^2\right)-\frac{1}{\nu}}
{1+y_3^2}.
\]

Therefore these expressions are non-positive for all \(x_3,y_3\in\mathbb R\)
if and only if
\[
K_{12}-e_3^2\leq 0,
\qquad
K_{23}\leq 0,
\qquad
K_{31}\leq 0.
\]
The last two inequalities are automatic, because
\[
K_{23}=K_{31}=-\frac{1}{\nu}<0.
\]
The first condition gives
\[
\frac{1}{\nu}-e_3^2\leq 0,
\]
or equivalently
\[
|e_3|\geq \frac{1}{\sqrt{\nu}}.
\]

\subsubsection{The group NIL}
%\subsection*{The group Nil}

The left-invariant metric on Nil is given by
\[
\begin{pmatrix}
\lambda & 0 & 0\\
0 & \lambda & 0\\
0 & 0 & 1
\end{pmatrix},
\qquad \lambda>0.
\]
We use the orthonormal Milnor base
\[
Y_1=\frac{1}{\sqrt{\lambda}}X,\qquad
Y_2=\frac{1}{\sqrt{\lambda}}Y,\qquad
Y_3=Z.
\]
with structural constants
\[
[Y_1,Y_2]=\frac{1}{\lambda}Y_3,\qquad
[Y_2,Y_3]=0,\qquad
[Y_3,Y_1]=0.
\]
Thus the only non-zero structure constant is
\[
a_3=: \frac{1}{\lambda}.
\]
The sectional curvatures are
\[
K(Y_1,Y_2)=-\frac{3}{4\lambda^2},
\qquad
K(Y_2,Y_3)=\frac{1}{4\lambda^2},
\qquad
K(Y_3,Y_1)=\frac{1}{4\lambda^2}.
\]

An automorphism of Nil is of form
\[
\begin{pmatrix}
a & c & 0\\
b & d & 0\\
* & * & ad-bc
\end{pmatrix},
\qquad ad-bc\neq 0.
\]
Therefore, with respect to the Milnor base, we may write
\[
\varphi(Y_1)=aY_1+bY_2+rY_3,
\]
\[
\varphi(Y_2)=cY_1+dY_2+sY_3,
\]
\[
\varphi(Y_3)=(ad-bc)Y_3.
\]
Assume that $\varphi(Y_i)$ are scaled such that $\varphi(Y_i) = 1$. This yields
$\|\varphi(Y_3)\|^2=(ad-bc)^2,$ and thus $(ad-bc)^2 = 1$. By restricting to the connected component of the identity, we may assume $ad - bc = 1$. \\~\\
We also assume that all the planes are spanned by orthogonal bases and $\varphi$ maps the base $\{Y_1,Y_2,Y_3\}$ to the orthogonal base $\{\varphi(Y_1),\;\varphi(Y_2),\; \varphi(Y_3)\}$. Since $
\langle \varphi(Y_1),\varphi(Y_3)\rangle=r(ad-bc)=r$ and $\langle \varphi(Y_2),\varphi(Y_3)\rangle=s(ad-bc)=s,$ we obtain that $r=s=0$ by orthogonality. \\~\\
Thus
\[
\varphi(Y_1)=aY_1+bY_2,
\qquad
\varphi(Y_2)=cY_1+dY_2,
\qquad
\varphi(Y_3)=Y_3.
\]
and the matrix
\[
\begin{pmatrix}
a & c\\
b & d
\end{pmatrix} \qquad \text{lies in SO(2).}
\]
Therefore we can write
\[
\varphi(Y_1)=\cos\theta\,Y_1+\sin\theta\,Y_2,
\]
\[
\varphi(Y_2)=-\sin\theta\,Y_1+\cos\theta\,Y_2,
\]
Let
\[
E=e_1Y_1+e_2Y_2+e_3Y_3.
\]
We now compute the Weyl sectional curvatures of the planes spanned by
\(\varphi(Y_i)\) and \(\varphi(Y_j)\). Since the vectors
\(\varphi(Y_1),\varphi(Y_2),\varphi(Y_3)\) are orthonormal, we can use that the Weyl sectional curvature is given by
\[\widehat K_E(X,Y)=K(X,Y)+\langle Y,E\rangle^2+\langle X,E\rangle^2-\langle E,E\rangle-\langle Y,\alpha(Y,E)\rangle-\langle \alpha(X,E),X\rangle.\]

First, we observe that the linear terms vanish. Indeed, for any horizontal
unit vector
\[
U=uY_1+vY_2
\]
one has
\[
\langle \alpha(U,E),U\rangle=0.
\]
Also,
\[
\langle \alpha(Y_3,E),Y_3\rangle=0.
\]
Hence the formula for the Weyl sectional curvature reduces to
\[\widehat K_E(X,Y)=K(X,Y)+\langle X,E\rangle^2+\langle Y,E\rangle^2-\langle E,E\rangle.\]
Equivalently, if \(N\) is the unit normal vector to the plane
\(\operatorname{span}(X,Y)\), then
\[
\widehat K_E(X,Y)
=
K(X,Y)-\langle E,N\rangle^2.
\]

For the first plane we have
\[
\operatorname{span}(\varphi(Y_1),\varphi(Y_2))
=
\operatorname{span}(Y_1,Y_2).
\]
Its unit normal vector is \(Y_3\). Therefore
\[
\widehat K_E(\varphi(Y_1),\varphi(Y_2))= K(Y_1,Y_2)-e_3^2 = -\frac{3}{4\lambda^2}-e_3^2.
\]
This  is always negative. For the plane, $\operatorname{span}(\varphi(Y_2),\varphi(Y_3)),$ the unit normal vector is \(\varphi(Y_1)\). Hence $\widehat K_E(\varphi(Y_2),\varphi(Y_3)) = K(Y_2,Y_3)-\langle E,\varphi(Y_1)\rangle^2.$ Since $\varphi(Y_1)=\cos\theta\,Y_1+\sin\theta\,Y_2,$ we obtain that $\langle E,\varphi(Y_1)\rangle=e_1\cos\theta+e_2\sin\theta$.
Therefore
\[
\widehat K_E(\varphi(Y_2),\varphi(Y_3))=\frac{1}{4\lambda^2}-(e_1\cos\theta+e_2\sin\theta)^2.
\]

For the third plane, $\operatorname{span}(\varphi(Y_3),\varphi(Y_1)),$
the unit normal vector is \(\varphi(Y_2)\). Hence $\widehat K_E(\varphi(Y_3),\varphi(Y_1))=K(Y_3,Y_1)-\langle E,\varphi(Y_2)\rangle^2.$
Since $\varphi(Y_2)=-\sin\theta\,Y_1+\cos\theta\,Y_2,$ we get that $\langle E,\varphi(Y_2)\rangle=-e_1\sin\theta+e_2\cos\theta$.
Therefore
\[
\widehat K_E(\varphi(Y_3),\varphi(Y_1))
=
\frac{1}{4\lambda^2}
-
(-e_1\sin\theta+e_2\cos\theta)^2.
\]

We now check whether these expressions can be non-positive for every
automorphic isometry \(\varphi\), or equivalently for every \(\theta\). \\~\\
Let $E_H=e_1Y_1+e_2Y_2$
be the horizontal part of \(E\). Assume \(E_H=0\), i.e.
\[
e_1=e_2=0.
\]
Therefore
\[
\widehat K_E(\varphi(Y_2),\varphi(Y_3))=\widehat K_E(\varphi(Y_3),\varphi(Y_1))=\frac{1}{4\lambda^2}>0.
\]
Therefore for $E_H \neq 0$, there is a positive Weyl sectional curvature. Assume \(E_H\neq 0\). Then we may choose \(\theta\) such that $\varphi(Y_1)\perp E_H$, i.e. $\langle E,\varphi(Y_1)\rangle=0.$ Consequently,
\[
\widehat K_E(\varphi(Y_2),\varphi(Y_3))
=
\frac{1}{4\lambda^2}>0.
\]
Thus also in this case at least one of the Weyl sectional curvatures is
positive.

Therefore, for every left-invariant vector field \(E\), there exists an
automorphic isometry \(\varphi\) such that
\[
\widehat K_E(\varphi(Y_i),\varphi(Y_j))>0
\]
for some \(i\neq j\). Hence Nil admits no non-positive left-invariant Weyl
connections.
\begin{thebibliography}{ABCD}

\bibitem{A} D. V. Alekseevskii, {\it Homogeneos Riemannian spaces of negative curvature}, Math. USSR Sbornik, 25(1975), 87-109.
\bibitem{BJT} M.Boche\'nski, P. Jastrz\c ebski, A. Tralle,  {\it Stretched non-positive Weyl connections on solvable Lie groups}, Ann.  Mat. Pura  Appl., 203(2024), 1463-1481
\bibitem{DGP} C. Draper, A. Garvin, F. J. Palomo, {\it Invariant affine connections on odd-dimensional spheres}, Ann. Glob. Anal. Geom., 49(2016), 213-251.

\bibitem{H} W. G. Hoover, {\it Molecular Dynamics}, Lect. Notes Phys. vol. 258, Springer, 1986. 

\bibitem{KN} S. Kobayashi, K. Nomizu, {\it Foundations of Differential Geometry}, vol. 2, Interscience, 1969.
\bibitem{KHL} Ku Yong Ha, Jong Bum Lee, {\it Left-invariant metrics and curvatures on simply connected three dimensional Lie groups,} Math. Nachr., 82(2009), 868-898
\bibitem{M} J. Milnor, {\it Curvature of left-invariant metrics on Lie groups,} Adv. Math., 21(1976), 293-329.
\bibitem{N} K. Nomizu, {\it Affine connections on homogeneous spaces}, Amer. J. Math., 76(1954), 33-65.

\bibitem{TW} G. Tereszkiewicz, M. P. Wojtkowski, {\it Homogeneous Weyl connections of non-positive curvature}, Ann. Global. Anal., 51(2017), 209-229.

\bibitem{W} M. P. Wojtkowski, {\it W-flows on Weyl manifolds and Gaussian theromostats.} J.  Math. Pure Appl., 10(2000),953-974

\bibitem{V} S. Van Tuong, {\it Metrics on 4-dimensional solvable unimodular Lie groups}, Ann. Global Anal. Geom., 51(2016), 109-128 
\end{thebibliography}
\end{document}
